% GENERATED FILE. Edit canonical lessons, then run npm run generate:books.
% canonical-source-hash: fnv1a64:48c483015ee692f7
% canonical-lessons: RU-C511-turnir, RU-C511-chempionat, RU-C511-chempion, RU-C511-medal, RU-C511-nagrazhdat

\chapter{Medals and Champions}
\label{ch:ru-medals-and-champions}
\hlchaptermodality{511}

\begin{chapteropening}
\textbf{By the end of this chapter:} \emph{I can talk about tournaments, champions and medals.}

Complete the last lesson of chapter 511: I can talk about tournaments, champions and medals.
\end{chapteropening}

\section[\texorpdfstring{\ru{турнир} (turnír)}{турнир (turnír)}]{\ru{турнир} (turnír) — a tournament}

\label{lesson:RU-C511-turnir}



\begin{quote}
Before the new one: say the Russian for a loss, then the Russian for to be in the lead.
\end{quote}

\subsection*{You'll want to know: \ru{турнир}}
\textbf{\ru{турнир}} (\emph{turnír}) — \textquotedblleft{}a tournament\textquotedblright{}.

\begin{quote}
\ru{В} \ru{турнире} \ru{участвуют} \ru{восемь} \ru{команд}.
\end{quote}

\begin{quote}
\emph{V turníre uchástvuyut vósem kománd.}
\end{quote}

\begin{quote}
Eight teams are taking part in the tournament.
\end{quote}

Masculine.

\begin{cousinweb}
EtymDB lists \emph{\ru{турнир}} without a source; it is usually traced through German \emph{Turnier} to Old French \emph{tornei}.
\end{cousinweb}

\begin{grammarlens}[title={What you've built}]
One more word for telling, explaining and describing at B1.
\end{grammarlens}

\subsection*{Your turn}
\begin{itemize}
\raggedright
  \item \emph{Say it:} \emph{turnír}
  \item \emph{Say it:} \emph{turnír}, once more
  \item \emph{Say it:} a sentence of your own with \emph{turnír}
  \item \emph{Recall:} say the Russian for a loss, then the Russian for to be in the lead, then say \emph{turnír} again
\end{itemize}

\begin{tcolorbox}[breakable,colback=teal!4,colframe=teal!35!black,title={Before you move on}]
What does \ru{турнир} mean? (\textquotedblleft{}A tournament\textquotedblright{}.) Say it once more.
\end{tcolorbox}
\section[\texorpdfstring{\ru{чемпионат} (chempionát)}{чемпионат (chempionát)}]{\ru{чемпионат} (chempionát) — a championship}

\label{lesson:RU-C511-chempionat}



\begin{quote}
Before the new one: say the Russian for to be in the lead, then the Russian for a tournament.
\end{quote}

\subsection*{You'll want to know: \ru{чемпионат}}
\textbf{\ru{чемпионат}} (\emph{chempionát}) — \textquotedblleft{}a championship\textquotedblright{}.

\begin{quote}
\ru{Чемпионат} \ru{мира} \ru{по} \ru{футболу} \ru{смотрят} \ru{миллионы} \ru{людей}.
\end{quote}

\begin{quote}
\emph{Chempionát míra po futbólu smótryat millióny lyudéy.}
\end{quote}

\begin{quote}
Millions of people watch the football World Cup.
\end{quote}

Masculine. \emph{\ru{Чемпионат} \ru{мира}} — a world championship.

\begin{cousinweb}
\emph{\ru{чемпионат}} is built on \emph{\ru{чемпион}} 'a champion', which was borrowed from English \emph{champion} (EtymDB).
\end{cousinweb}

\begin{grammarlens}[title={What you've built}]
One more word for telling, explaining and describing at B1.
\end{grammarlens}

\subsection*{Your turn}
\begin{itemize}
\raggedright
  \item \emph{Say it:} \emph{chempionát}
  \item \emph{Say it:} \emph{chempionát}, once more
  \item \emph{Say it:} a sentence of your own with \emph{chempionát}
  \item \emph{Recall:} say the Russian for to be in the lead, then the Russian for a tournament, then say \emph{chempionát} again
\end{itemize}

\begin{tcolorbox}[breakable,colback=teal!4,colframe=teal!35!black,title={Before you move on}]
What does \ru{чемпионат} mean? (\textquotedblleft{}A championship\textquotedblright{}.) Say it once more.
\end{tcolorbox}
\section[\texorpdfstring{\ru{чемпион} (chempión)}{чемпион (chempión)}]{\ru{чемпион} (chempión) — a champion}

\label{lesson:RU-C511-chempion}



\begin{quote}
Before the new one: say the Russian for a tournament, then the Russian for a championship.
\end{quote}

\subsection*{You'll want to know: \ru{чемпион}}
\textbf{\ru{чемпион}} (\emph{chempión}) — \textquotedblleft{}a champion\textquotedblright{}.

\begin{quote}
\ru{Он} \ru{уже} \ru{три} \ru{года} \ru{чемпион} \ru{страны}.
\end{quote}

\begin{quote}
\emph{On uzhé tri góda chempión straný.}
\end{quote}

\begin{quote}
He has been the national champion for three years.
\end{quote}

Masculine; feminine \ru{чемпионка}.

\begin{cousinweb}
\emph{\ru{чемпион}} was borrowed from English \emph{champion} (EtymDB).
\end{cousinweb}

\begin{grammarlens}[title={What you've built}]
One more word for telling, explaining and describing at B1.
\end{grammarlens}

\subsection*{Your turn}
\begin{itemize}
\raggedright
  \item \emph{Say it:} \emph{chempión}
  \item \emph{Say it:} \emph{chempión}, once more
  \item \emph{Say it:} a sentence of your own with \emph{chempión}
  \item \emph{Recall:} say the Russian for a tournament, then the Russian for a championship, then say \emph{chempión} again
\end{itemize}

\begin{tcolorbox}[breakable,colback=teal!4,colframe=teal!35!black,title={Before you move on}]
What does \ru{чемпион} mean? (\textquotedblleft{}A champion\textquotedblright{}.) Say it once more.
\end{tcolorbox}
\section[\texorpdfstring{\ru{медаль} (medál)}{медаль (medál)}]{\ru{медаль} (medál) — a medal}

\label{lesson:RU-C511-medal}



\begin{quote}
Before the new one: say the Russian for a championship, then the Russian for a champion.
\end{quote}

\subsection*{You'll want to know: \ru{медаль}}
\textbf{\ru{медаль}} (\emph{medál}) — \textquotedblleft{}a medal\textquotedblright{}.

\begin{quote}
\ru{Спортсмен} \ru{получил} \ru{золотую} \ru{медаль}.
\end{quote}

\begin{quote}
\emph{Sportsmén poluchíl zolotúyu medál.}
\end{quote}

\begin{quote}
The athlete won a gold medal.
\end{quote}

Feminine.

\begin{cousinweb}
EtymDB lists \emph{\ru{медаль}}, glossed 'medal', without a source; it is usually traced through French \emph{médaille} to Italian \emph{medaglia}.
\end{cousinweb}

\begin{grammarlens}[title={What you've built}]
One more word for telling, explaining and describing at B1.
\end{grammarlens}

\subsection*{Your turn}
\begin{itemize}
\raggedright
  \item \emph{Say it:} \emph{medál}
  \item \emph{Say it:} \emph{medál}, once more
  \item \emph{Say it:} a sentence of your own with \emph{medál}
  \item \emph{Recall:} say the Russian for a championship, then the Russian for a champion, then say \emph{medál} again
\end{itemize}

\begin{tcolorbox}[breakable,colback=teal!4,colframe=teal!35!black,title={Before you move on}]
What does \ru{медаль} mean? (\textquotedblleft{}A medal\textquotedblright{}.) Say it once more.
\end{tcolorbox}
\section[\texorpdfstring{\ru{награждать} (nagrazhdát)}{награждать (nagrazhdát)}]{\ru{награждать} (nagrazhdát) — to award}

\label{lesson:RU-C511-nagrazhdat}



\begin{quote}
Before the new one: say the Russian for a champion, then the Russian for a medal.
\end{quote}

\subsection*{You'll want to know: \ru{награждать}}
\textbf{\ru{награждать}} (\emph{nagrazhdát}) — \textquotedblleft{}to award\textquotedblright{}.

\begin{quote}
\ru{Победителей} \ru{награждают} \ru{медалями}.
\end{quote}

\begin{quote}
\emph{Pobedíteley nagrazhdáyut medályami.}
\end{quote}

\begin{quote}
The winners are awarded medals.
\end{quote}

Imperfective; perfective \ru{наградить} (\ru{награжу}). \emph{\ru{Награждать} \ru{кого} \ru{чем}}.

\begin{cousinweb}
\emph{\ru{награждать}} is built on \emph{\ru{награда}} 'a reward'; EtymDB has no entry, and the origin of \emph{\ru{награда}} is not settled.
\end{cousinweb}

\begin{grammarlens}[title={What you've built}]
One more word for telling, explaining and describing at B1.
\end{grammarlens}

\subsection*{Your turn}
\begin{itemize}
\raggedright
  \item \emph{Say it:} \emph{nagrazhdát}
  \item \emph{Say it:} \emph{nagrazhdát}, once more
  \item \emph{Say it:} a sentence of your own with \emph{nagrazhdát}
  \item \emph{Recall:} say the Russian for a champion, then the Russian for a medal, then say \emph{nagrazhdát} again
\end{itemize}

\begin{tcolorbox}[breakable,colback=teal!4,colframe=teal!35!black,title={Before you move on}]
What does \ru{награждать} mean? (\textquotedblleft{}To award\textquotedblright{}.) Say it once more.
\end{tcolorbox}
